\chapter{BX-32: Embedded Floating Point with Zve32f}
\label{ch:bx32}

\section{The first ratified claim}

BX-32 is an RV32 hart whose data path is thirty-two bits, with the scalar extensions you need for a hosted compiler, and with Zve32f. This is the first machine in the book that is allowed to tell a standard toolchain ``vector, with single-precision float''. The discovery path is the standard one. The custom CSR used by SVM-8 and SVM-16 may remain, returning a code that means ``superset of SVM-16'', but \texttt{misa} and the vector CSRs are authoritative.

\begin{spec}
ELEN = 32. VLEN is at least 128.\\
Legal SEW = 8, 16, and 32.\\
FP32 vector arithmetic, conversions, and fused multiply-add.\\
Indexed loads and stores.\\
Integer behaviour of SVM-16 remains legal.
\end{spec}

A loop assembled for SVM-8, using only EEW=8, runs. A loop assembled for SVM-16, using EEW=16 and \texttt{vwmul}, runs. A loop that sets SEW=32 runs for the first time.

\section{Floating point in the lane}

Each lane that supports FP32 contains a single-precision adder and multiplier, or a fused unit if the FPGA's DSP blocks map cleanly onto one. Fused multiply-add is required by the extension. Rounding follows the dynamic rounding mode of the hart.

Do not share one FP unit across all lanes by silent time-multiplexing unless the cycle count is in the profile card. A programmer who reads Zve32f will expect four lanes to retire four FP32 elements per cycle once the pipeline is full. If your first image has one FP unit and four integer lanes, say ``integer lanes: 4, FP lanes: 1'' on the card, and do not claim that the peak is four times the FP unit.

\section{Where SVM-16 software feels the change}

The change a programmer notices is not the wider bus. It is that SEW=32 no longer traps, and that a matrix kernel can stay in floating point instead of in widened integers. Xsvm, if it is still elaborated, must use the same SEW as the surrounding code. A four-by-four of FP32 values in \texttt{v8}--\texttt{v11} is the useful form. The integer form remains for SVM-16 compatibility.

Prefer, in new BX-32 software, the standard sequence of fused multiply-adds over Xsvm. Keep Xsvm for kernels you have measured, and for comparison with the BX-16 geometry image.

\section{Two harts}

A second RV32 hart is optional. It shares memory through an arbiter you already know how to write, because BX-16 has several masters. The second hart is for a concurrent system, not a different instruction set. Both harts, if both are present, implement Zve32f or the second is a scalar only and the profile card says so. Do not build one hart with FP32 vector and another with a private opcode map.

Local scratchpads are ordinary memories in the address map, marked cacheable or not according to your platform. They are not a third address space with a private load instruction. Vector loads to them use the same unit-stride encoding as vector loads to the rest of memory.
